\documentclass[aspectratio=169]{beamer}
\usepackage{nagoyatheme}

\usepackage[T1]{fontenc}
\usepackage[dvipsnames]{xcolor}
\usepackage{amsmath, amssymb}
\usepackage{mathpartir}  % \inferrule

\usepackage{tikz}
\usetikzlibrary{arrows.meta, positioning}

% Rocq code rendered by Alectryon (see Makefile)
\usepackage{alectryon}
\usepackage{pygments}
\usepackage{appendixnumberbeamer}
\DeclareUnicodeCharacter{0393}{\ensuremath{\Gamma}}
\DeclareUnicodeCharacter{0394}{\ensuremath{\Delta}}
\DeclareUnicodeCharacter{22A2}{\ensuremath{\vdash}}
% smaller Alectryon goal boxes on dense slides (default is \small)
\makeatletter
\newcommand\alectryonscriptoutput{\renewcommand*{\alectryon@output@setsize}{\scriptsize}}
\makeatother
% VS Code-like proof state for step-by-step proofs: numbered goals,
% no extra space above the output (top aligned with the code).
\makeatletter
\newcounter{alecgoal}
\newcommand\alectryonvscode{%
  \renewcommand*{\alectryon@output@setsize}{\scriptsize}%
  \renewenvironment{alectryon@env}{\par\nointerlineskip}{\par}%
  \renewenvironment{alectryon@goals}
    {\setcounter{alecgoal}{0}\begin{alectryon@inner}%
     \def\alectryon@sep{\alectryon@unskip\ignorespaces}\ignorespaces}
    {\alectryon@unskip\end{alectryon@inner}}%
  \renewenvironment{alectryon@goal}
    {\stepcounter{alecgoal}%
     \ifnum\value{alecgoal}>1 \renewcommand\alectryon@infrule[1]{}%
       \def\alectryon@outer@vsep{}\par\vspace{0.6ex}\fi
     \def\alectryon@sep{\alectryon@unskip\alectryon@goal@sep\ignorespaces}\ignorespaces}
    {\alectryon@unskip\par}%
  \renewenvironment{alectryon@conclusion}
    {\par\leavevmode\strut\arabic{alecgoal}\ :\ \ignorespaces}{\strut\par}}
% a Rocq message that Alectryon does not capture, in the same style
\newcommand\rocqmsg[1]{\begin{tcolorbox}[colback=alectryon@tango@light@aluminium,
    colframe=alectryon@tango@medium@aluminium, boxrule=2.5pt, arc=0pt,
    left=2pt, right=2pt, top=2pt, bottom=2pt, boxsep=0pt]\ttfamily\scriptsize #1\end{tcolorbox}}
\makeatother
% step-by-step proofs: \mk{k} = cursor shown on overlay k (0 = never)
\newcommand\mk[1]{\makebox[1.1em][l]{\ifnum#1>0 \alt<#1>{\textcolor{nagoyaGreen}{\textbf{>}}}{}\fi}}
% \popup<k>{...}: takeaway box over the slide on overlay k
\usetikzlibrary{shadows}
\newcommand<>\popup[1]{\only#2{\begin{tikzpicture}[remember picture,overlay]
  \node[fill=white, draw=nagoyaGreen, very thick, rounded corners, drop shadow,
        inner sep=10pt] at (current page.center)
        {\begin{minipage}{.7\paperwidth}\small #1\end{minipage}};
\end{tikzpicture}}}

\definecolor{ColorOne}{named}{nagoyaGreen}
\definecolor{ColorTwo}{named}{Gray}
\definecolor{ColorThree}{named}{Plum}

\setbeamerfont{title}{size=\Large,series=\bfseries}  % fit the title on two lines
\title[MLTT in Rocq with mutual inductive types]{Implementing MLTT in Rocq with mutual inductive types:\\ challenges and techniques}
\author[WULFMAN]{Pierre-Emmanuel WULFMAN}
\institute[Nagoya University]{Graduate School of Mathematics, Nagoya University}
\date{TPP 2026 --- RIKEN AIP, Tokyo --- 16 November 2026}

\begin{document}

% =====================================================================
% Part 1 -- The problem
% =====================================================================

% --- Slide 1 ---------------------------------------------------------
\begin{frame}[plain, noframenumbering]
  \titlepage
\end{frame}

% --- Slide 2 ---------------------------------------------------------
\section{The problem}
\begin{frame}{Context}
  \begin{itemize}
    \item Implementing in Rocq the extensional MLTT of Angiuli \& Gratzer,
          \emph{Principles of Dependent Type Theory}
    \item A dependent type theory is a set of \textbf{cyclic definitions}:
          contexts, types, terms and substitutions all refer to each other
    \item This talk: what I learned \emph{as a Rocq user} about
          \begin{itemize}
            \item defining cyclic definitions
            \item stating and proving lemmas about them
          \end{itemize}
    \item The lessons apply whenever your definitions are cyclic
    \item Extensional MLTT is particularly relevant: 4 judgements and their 4 equality judgements,
          \textbf{8 mutually inductive} judgements, far more interwoven than most definitions
  \end{itemize}
\end{frame}

% --- Slide 3 ---------------------------------------------------------
\begin{frame}{Simple types: no cycle}
  Simply typed $\lambda$-calculus: a definition order exists.
  \medskip

  \begin{columns}[c]
    \column{.5\textwidth}
    \small
    \[
      \begin{array}{rcl}
        A & ::= & \mathsf{Bool} \mid A \to A\\
        t & ::= & x \mid \lambda x{:}A.\,t \mid t\ t \mid \dots\\
        \Gamma & ::= & \cdot \mid \Gamma, x{:}A
      \end{array}
    \]
    \[
      \inferrule{\Gamma(x) = A}{\Gamma \vdash x : A}
    \qquad
      \inferrule{\Gamma, x{:}A \vdash t : B}{\Gamma \vdash \lambda x{:}A.\,t : A \to B}
    \]
    \[
      \inferrule{\Gamma \vdash t_1 : A \to B \\ \Gamma \vdash t_2 : A}{\Gamma \vdash t_1\ t_2 : B}
    \]

    \column{.45\textwidth}
    \centering
    \begin{tikzpicture}[>=Stealth, every node/.style={font=\small, draw=ColorOne, rounded corners, inner sep=4pt},
                        dep/.style={->, thick, ColorOne}]
      \node (jd) at (1.3,2.4) {$\Gamma \vdash t : A$};
      \node (cx) at (0,1.2)   {contexts $\Gamma$};
      \node (tm) at (2.6,1.2) {terms $t$};
      \node (ty) at (1.3,0)   {types $A$};
      \draw[dep] (jd) -- (cx);
      \draw[dep] (jd) -- (tm);
      \draw[dep] (jd) -- (ty);
      \draw[dep] (cx) -- (ty);
      \draw[dep] (tm) -- (ty);
    \end{tikzpicture}

    \smallskip
    {\footnotesize $X \to Y$: $X$ refers to $Y$}
  \end{columns}

  \medskip
  \centering
  \textbf{No cycle: define types, then terms and contexts, then typing.}

  \vfill
  {\scriptsize Rules as in Software Foundations, Vol.~2 \emph{Programming Language Foundations}, chapter \texttt{Stlc}.}
\end{frame}

% --- Slide 4 ---------------------------------------------------------
\begin{frame}{Dependent types: a cycle}
  Dependent type theory: the judgements need each other (a few of the rules).
  \medskip

  \begin{columns}[c]
    \column{.57\textwidth}
    \centering\small
    \setlength{\tabcolsep}{3pt}
    \begin{tabular}{cc}
      $\inferrule{\vdash \Gamma \ \mathsf{cx} \\ \Gamma \vdash A \ \mathsf{type}}
                 {\vdash \Gamma . A \ \mathsf{cx}}$
      &
      $\inferrule{\Delta \vdash \gamma : \Gamma \\ \Gamma \vdash A \ \mathsf{type}}
                 {\Delta \vdash A[\gamma] \ \mathsf{type}}$
      \\[4ex]
      $\inferrule{\Delta \vdash \gamma : \Gamma \\\\ \Gamma \vdash A \ \mathsf{type}
                  \\ \Delta \vdash a : A[\gamma]}
                 {\Delta \vdash \gamma . a : \Gamma . A}$
      &
      $\inferrule{\Delta \vdash \gamma : \Gamma \\ \Gamma \vdash a : A}
                 {\Delta \vdash a[\gamma] : A[\gamma]}$
    \end{tabular}

    \smallskip
    {\footnotesize \dots\ and many more rules}

    \column{.41\textwidth}
    \centering
    \begin{tikzpicture}[>=Stealth, every node/.style={font=\small, draw=ColorOne, rounded corners, inner sep=4pt},
                        dep/.style={->, thick, ColorOne}]
      \node (cx) at (0,2)   {$\vdash \Gamma\ \mathsf{cx}$};
      \node (ty) at (2.4,2) {$\Gamma \vdash A\ \mathsf{type}$};
      \node (sb) at (0,0)   {$\Delta \vdash \gamma : \Gamma$};
      \node (tm) at (2.4,0) {$\Gamma \vdash a : A$};
      \draw[dep] (cx) to[bend left=15] (ty);
      \draw[dep] (ty) to[bend left=15] (cx);
      \draw[dep] (ty) to[bend left=15] (sb);
      \draw[dep] (sb) to[bend left=15] (ty);
      \draw[dep] (sb) to[bend left=15] (tm);
      \draw[dep] (tm) to[bend left=15] (sb);
      \draw[dep] (sb) -- (cx);
      \draw[dep] (tm) -- (ty);
    \end{tikzpicture}

    \smallskip
    {\footnotesize $X \to Y$: $X$ refers to $Y$\\
     + 4 equality judgements, same cycle}
  \end{columns}

  \medskip
  \centering
  \textbf{No definition order: everything at once.}

  \vfill
  {\scriptsize Rules from Angiuli \& Gratzer, \emph{Principles of Dependent Type Theory}, Sec.~2.3.\\
   In the code: $\Gamma.A$ is written \texttt{Γ \# A} (a ``\texttt{.}'' would end the Rocq sentence), and $\gamma$ is written \texttt{y}.}
\end{frame}

% =====================================================================
% Part 2 -- Two designs
% =====================================================================

% --- [M1-def] ------------------------------------------------------
\section{Two designs}
\begin{frame}{Method 1: one family indexed by a sort}
  The usual advice: one inductive, indexed by which judgement it is.
\begin{alectryon}{coq}{none}
Inductive preSubs : Type :=
| psub_id
| psub_weak
| psub_ext (y : preSubs) (a : preTerm)
with preTerm : Type :=
| ptrm_qar
| ptrm_subs (a : preTerm) (y : preSubs).
Inductive preType : Type :=
| ptyp_subs (A : preType) (y : preSubs)
| ptyp_unit.
Inductive preCtx : Type :=
| pctx_one
| pctx_ext (Γ : preCtx) (A : preType).
Declare Scope ctx_scope. Delimit Scope ctx_scope with ctx.
#[add_top] Bind Scope ctx_scope with preCtx.
Declare Scope subs_scope. Delimit Scope subs_scope with subs.
#[add_top] Bind Scope subs_scope with preSubs.
Declare Scope ty_scope. Delimit Scope ty_scope with ty.
#[add_top] Bind Scope ty_scope with preType.
Declare Scope term_scope. Delimit Scope term_scope with term.
#[add_top] Bind Scope term_scope with preTerm.
Notation "Γ # A" := (pctx_ext Γ A) (at level 2, left associativity) : ctx_scope.
Notation "y # a" := (psub_ext y a) (at level 2, left associativity) : subs_scope.
Notation "A [ y ]" := (ptyp_subs A y) (at level 1) : ty_scope.
Notation "a [ y ]" := (ptrm_subs a y) (at level 1) : term_scope.
Abbreviation Unit := ptyp_unit.
Abbreviation p := psub_weak.

\end{alectryon}
  \vspace{-1ex}
  \medskip
  \small
\begin{alectryon}{coq}{none}
Module Method1.
\end{alectryon}
\begin{alectryon}{coq}{in}
Inductive Sort :=
| CX (Γ : preCtx)
| TY (Γ : preCtx) (A : preType)
| TM (Γ : preCtx) (A : preType) (a : preTerm)
| SB (Δ Γ : preCtx) (y : preSubs).
Inductive J : Sort -> Prop :=
| ctx_ext Γ A : J (CX Γ) -> J (TY Γ A) -> J (CX (Γ # A))
| type_subs Δ Γ y A : J (SB Δ Γ y) -> J (TY Γ A) -> J (TY Δ A[y])
| term_subs Δ Γ y A a : J (SB Δ Γ y) -> J (TM Γ A a) -> J (TM Δ A[y] a[y])
| subs_weak Γ A : J (TY Γ A) -> J (SB (Γ # A) Γ p)
| subs_ext Δ Γ y A a :
    J (SB Δ Γ y) -> J (TY Γ A) -> J (TM Δ A[y] a)
    -> J (SB Δ (Γ # A) (y # a))
(* ... and many more rules *).
\end{alectryon}
\begin{alectryon}{coq}{none}
End Method1.
\end{alectryon}
\end{frame}

% --- [M2-def] ------------------------------------------------------
\begin{frame}{Method 2: one family per judgement, tied with \texttt{with}}
  What I did instead: the natural step from STLC.
  \small
\begin{alectryon}{coq}{none}
Module Method2.
\end{alectryon}
\begin{alectryon}{coq}{in}
Inductive CtxJG : preCtx -> Prop :=
| ctx_ext Γ A : CtxJG Γ -> TypeJG Γ A -> CtxJG (Γ # A)
with TypeJG : preCtx -> preType -> Prop :=
| type_subs Δ Γ y A : SubsJG Δ Γ y -> TypeJG Γ A -> TypeJG Δ A[y]
with TermJG : preCtx -> preType -> preTerm -> Prop :=
| term_subs Δ Γ y A a : SubsJG Δ Γ y -> TermJG Γ A a -> TermJG Δ A[y] a[y]
with SubsJG : preCtx -> preCtx -> preSubs -> Prop :=
| subs_weak Γ A : TypeJG Γ A -> SubsJG (Γ # A) Γ p
| subs_ext Δ Γ y A a :
    SubsJG Δ Γ y -> TypeJG Γ A -> TermJG Δ A[y] a
    -> SubsJG Δ (Γ # A) (y # a)
(* ... and many more rules, for each judgement *).
\end{alectryon}
\begin{alectryon}{coq}{none}
End Method2.
\end{alectryon}
\end{frame}

% =====================================================================
% Part 3 -- Proving lemmas
% =====================================================================

% --- [M2-ind] ------------------------------------------------------
\section{Proving lemmas}
\begin{frame}{Method 2: induction on a judgement}
  A well-formed type lives in a well-formed context.
\begin{alectryon}{coq}{none}
From Stdlib Require Import Program.Equality.
Module Ind2.
Import Method2.
Set Printing Width 75.
\end{alectryon}
  \small
\begin{alectryon}{coq}{}
Lemma wf Γ A : TypeJG Γ A -> CtxJG Γ. (* .in *)
Proof. (* .none *)
intros H; dependent induction H.
\end{alectryon}
\begin{alectryon}{coq}{none}
Abort.
End Ind2.
\end{alectryon}

  \vfill
  \begin{itemize}
    \item \texttt{IHTypeJG : CtxJG Γ} does not help: the goal is about \texttt{Δ}
    \item What is missing is visible: \texttt{H -> goal}, i.e.\ \texttt{SubsJG Δ Γ y -> CtxJG Δ}
    \item But proving \texttt{SubsJG Δ Γ y -> CtxJG Δ} will require \texttt{TypeJG Γ A -> CtxJG Γ}
  \end{itemize}
  \medskip
  \centering\textbf{We need to prove both lemmas together.}
\end{frame}

% --- [Scheme] ------------------------------------------------------
\begin{frame}{The manual's answer: \texttt{Scheme} / \texttt{Combined Scheme}}
  \small
  The manual's section on mutual inductive types points to \texttt{Scheme}.
  \medskip

\begin{alectryon}{coq}{none}
Module Sch.
Import Method2.
\end{alectryon}
  \begin{columns}[t,totalwidth=1.06\textwidth]
    \column{.56\textwidth}
    \scriptsize\alectryonvscode
\begin{alectryon}{coq}{in}
Scheme cx' := Minimality for CtxJG Sort Prop
with   ty' := Minimality for TypeJG Sort Prop
with   tm' := Minimality for TermJG Sort Prop
with   sb' := Minimality for SubsJG Sort Prop.
Combined Scheme JG_ind from cx', ty', tm', sb'.
Check JG_ind.
\end{alectryon}
\begin{alectryon}{coq}{none}
End Sch.
\end{alectryon}

    \column{.5\textwidth}
    \colorbox{gray!10}{\parbox[t]{.95\linewidth}{\ttfamily\scriptsize
      forall (P : preCtx -> Prop) \dots{} P2,\\
      \ \dots{} -> (forall Δ Γ y A,\\
      \ \ \ SubsJG Δ Γ y -> P2 Δ Γ y ->\\
      \ \ \ TypeJG Γ A -> P0 Γ A -> P0 Δ A[y]) -> \dots\\
      \ (forall p, CtxJG p -> P p) /\textbackslash{} \dots}}
    {\tiny (output of \texttt{Check} abridged; full type checked by Rocq 9.3)}
  \end{columns}

  \vfill
  \begin{itemize}
    \item one predicate per judgement (\texttt{P}, \texttt{P0}, \texttt{P1}, \texttt{P2}), one case per rule
    \item conclusion: one conjunct per judgement,
          \texttt{(forall p, CtxJG p -> P p) /\textbackslash{} (forall p p0, TypeJG p p0 -> P0 p p0) /\textbackslash{} \dots}
    \item my lemmas are separate statements: making them fit means a lot of boilerplate,
          so I looked for a more natural solution
  \end{itemize}
\end{frame}

% --- [Tried] -------------------------------------------------------
\begin{frame}{What I tried next}
  \begin{enumerate}
    \item[(a)] \textbf{Follow the cycle}: a second induction on the \texttt{SubsJG} premise, until it loops back
          \begin{itemize}
            \item the proof for \texttt{SubsJG} is written again inside the proof for \texttt{TypeJG}
            \item no single source of truth; with more lemmas, one copy per dependency
          \end{itemize}
    \item[(b)] \textbf{Break the cycle}: assume the other lemma
          \begin{itemize}
            \item prove the \texttt{TypeJG} lemma \emph{assuming} the \texttt{SubsJG} lemma
            \item prove the \texttt{SubsJG} lemma, using the first one to break the cycle
            \item restate the \texttt{TypeJG} lemma without the assumption
            \item correct, but boilerplate for every dependency
          \end{itemize}
    \item[(c)] \textbf{Combine the statements} with \texttt{/\textbackslash}:
          \texttt{(forall Γ A, TypeJG Γ A -> CtxJG Γ) /\textbackslash{} (forall Δ Γ y, \dots)}
          \begin{itemize}
            \item one common induction (e.g.\ on the context) needs a lot of extra structure
            \item a mutual principle avoids all that
          \end{itemize}
  \end{enumerate}
\end{frame}

% --- [Mutual] ------------------------------------------------------
\section{Mutual lemmas}
\begin{frame}[t]{(d) Mutual lemmas: \texttt{Lemma \dots\ with \dots}}
\begin{alectryon}{coq}{none}
Module Mut.
Import Method2.
Set Printing Width 48.
Lemma type_wf Γ A (H : TypeJG Γ A) : CtxJG Γ
with  subs_wf Δ Γ y (H : SubsJG Δ Γ y) :
        CtxJG Δ /\ CtxJG Γ.
\end{alectryon}
  \vspace{1.5ex}
  \begin{columns}[t,totalwidth=1.06\textwidth]
    \column{0.55\textwidth}
    {\ttfamily\scriptsize\setlength{\parindent}{0pt}
    \mbox{}\\[-1pt]
    \mk{0}\PY{k+kn}{Lemma}~\PY{n+nf}{type\PYZus{}wf}~\PY{n+nv}{Γ}~\PY{n+nv}{A}~\PY{o}{(}\PY{n+nv}{H}~\PY{o}{:}~\PY{n}{TypeJG}~\PY{n}{Γ}~\PY{n}{A}\PY{o}{)}~\PY{o}{:}~\PY{n}{CtxJG}~\PY{n}{Γ}\\
    \mk{0}\PY{k}{with}~~\PY{n}{subs\PYZus{}wf}~\PY{n}{Δ}~\PY{n}{Γ}~\PY{n}{y}~\PY{o}{(}\PY{n}{H}~\PY{o}{:}~\PY{n}{SubsJG}~\PY{n}{Δ}~\PY{n}{Γ}~\PY{n}{y}\PY{o}{)}~\PY{o}{:}\\
    \mk{0}~~~~~~~~\PY{n}{CtxJG}~\PY{n}{Δ}~\PY{o}{/}\PY{o}{\PYZbs{}}~\PY{n}{CtxJG}~\PY{n}{Γ}\PY{o}{.}\\
    \mk{1}\PY{k+kn}{Proof}\PY{o}{.}\\
    \mk{2}\PY{o}{\PYZhy{}}~\PY{n+nb}{dependent~induction}~\PY{n}{H}\PY{o}{.}\\
    \mk{3}~~\PY{o}{+}~\PY{n+nb+bp}{exact}~\PY{o}{(}\PY{n}{proj1}~\PY{o}{(}\PY{n}{subs\PYZus{}wf}~\PY{n}{\PYZus{}}~\PY{n}{\PYZus{}}~\PY{n}{\PYZus{}}~\PY{n}{H}\PY{o}{)}\PY{o}{)}\PY{o}{.}\\
    \mk{0}~~\PY{o}{+}~\PY{c}{(*}\PY{c}{~abridged~}\PY{c}{*)}\\
    \mk{4}\PY{o}{\PYZhy{}}~\PY{n+nb}{dependent~induction}~\PY{n}{H}\PY{o}{.}\\
    \mk{5}~~\PY{o}{+}~\PY{n+nb}{pose~proof}~\PY{o}{(}\PY{n}{type\PYZus{}wf}~\PY{n}{\PYZus{}}~\PY{n}{\PYZus{}}~\PY{n}{H}\PY{o}{)}~\PY{k}{as}~\PY{n}{HG}\PY{o}{.}\\
    \mk{6}~~~~\PY{n+nb}{split}\PY{o}{;}~\PY{o}{[}\PY{n+nb+bp}{exact}~\PY{o}{(}\PY{n}{ctx\PYZus{}ext}~\PY{n}{\PYZus{}}~\PY{n}{\PYZus{}}~\PY{n}{HG}~\PY{n}{H}\PY{o}{)}~\PY{o}{|}~\PY{n+nb+bp}{exact}~\PY{n}{HG}\PY{o}{]}\PY{o}{.}\\
    \mk{0}~~\PY{o}{+}~\PY{c}{(*}\PY{c}{~abridged~}\PY{c}{*)}\\
    \mk{7}\PY{k+kn}{Qed}\PY{o}{.}\\
    }
    \column{0.51\textwidth}
    \alectryonvscode\vspace*{-5pt}
    \begin{onlyenv}<1>
\begin{alectryon}{coq}{}
Proof. (* .unfold .no-in -.g#2.h{*} *)
\end{alectryon}
    \end{onlyenv}
    \begin{onlyenv}<2>
\begin{alectryon}{coq}{}
- (* .none *) dependent induction H. (* .unfold .no-in *)
\end{alectryon}
    \end{onlyenv}
    \begin{onlyenv}<3>
\begin{alectryon}{coq}{}
+ (* .none *) exact (proj1 (subs_wf _ _ _ H)). (* .none *)
\end{alectryon}
    \rocqmsg{This subproof is complete, but there are some unfocused goals.\\ Focus next goal with bullet +.}
    \end{onlyenv}
    \begin{onlyenv}<4>
\begin{alectryon}{coq}{}
- (* .none *) dependent induction H. (* .unfold .no-in -.g#2.h{*} *)
\end{alectryon}
    \end{onlyenv}
    \begin{onlyenv}<5>
\begin{alectryon}{coq}{}
+ (* .none *) pose proof (type_wf _ _ H) as HG. (* .unfold .no-in *)
\end{alectryon}
    \end{onlyenv}
    \begin{onlyenv}<6>
\begin{alectryon}{coq}{}
split; [exact (ctx_ext _ _ HG H) | exact HG]. (* .none *)
\end{alectryon}
    \rocqmsg{This subproof is complete, but there are some unfocused goals.\\ Focus next goal with bullet +.}
    \end{onlyenv}
\begin{alectryon}{coq}{none}
+ destruct IHSubsJG as [HD HG].
  split; [exact HD | exact (ctx_ext _ _ HG H0)].
\end{alectryon}
    \begin{onlyenv}<7>
\begin{alectryon}{coq}{}
Qed. (* .none *)
\end{alectryon}
    {\color{nagoyaGreen}\small$\checkmark$ accepted}
    \end{onlyenv}
  \end{columns}
\begin{alectryon}{coq}{none}

End Mut.
\end{alectryon}
  \popup<8>{\begin{itemize}
      \item combine only the lemmas that depend on each other
      \item they are all \textbf{in the context}, ready to be used
      \item scales well: no duplication, no boilerplate
    \end{itemize}}
\end{frame}

% --- [Level] -------------------------------------------------------
\begin{frame}[fragile]{Another level}
  \small
  The lemmas follow the shape of the judgements: one lemma per judgement of the cycle.
  In my development ({\footnotesize\texttt{2.3 - Calculus of Substitution/definitions.v}}):
\begin{semiverbatim}\scriptsize
Lemma subs_press \{Δ Γ : preCtx\} \{y : preSubs\} (H: [ Δ ⊢ y ~ Γ ]) :
        [ ⊢ Δ cx ] /\textbackslash [ ⊢ Γ cx ]
with type_press \{Γ : preCtx\} \{A : preType\} (H: [ Γ ⊢ A type ]) :
        [ ⊢ Γ cx ]
with eq_ctx_press \{Γ Δ : preCtx\} (H: eq_ctx Γ Δ) :
        [ ⊢ Γ cx ] /\textbackslash [ ⊢ Δ cx ]
with eq_type_press \{Γ : preCtx\} \{A B : preType\} (H: [ Γ ⊢ A == B type ]) :
        [ Γ ⊢ A type ] /\textbackslash [ Γ ⊢ B type ].
\end{semiverbatim}

  \begin{itemize}
    \item in each proof, the induction hypothesis is the \textbf{whole group}, not just the current lemma
    \item compared with (a)--(c) and \texttt{Scheme}: each lemma keeps its own statement;
          no motive to write, no conjunction to set up
    \item mentioned once in the manual; its section on mutual types points to \texttt{Scheme}
  \end{itemize}
\end{frame}

% --- [Power] -------------------------------------------------------
\begin{frame}[t]{With great power comes great responsibility}
  A circular ``proof'' is accepted, step by step\dots
\begin{alectryon}{coq}{none}
Module Bad.
Import Method2.
Set Printing Width 48.
Lemma type_wf Γ A (H : TypeJG Γ A) : CtxJG Γ
with  subs_wf Δ Γ y (H : SubsJG Δ Γ y) :
        CtxJG Δ /\ CtxJG Γ.
\end{alectryon}
  \vspace{1.5ex}
  \begin{columns}[t,totalwidth=1.06\textwidth]
    \column{0.55\textwidth}
    {\ttfamily\scriptsize\setlength{\parindent}{0pt}
    \mbox{}\\[-1pt]
    \mk{0}\PY{k+kn}{Lemma}~\PY{n+nf}{type\PYZus{}wf}~\PY{n+nv}{Γ}~\PY{n+nv}{A}~\PY{o}{(}\PY{n+nv}{H}~\PY{o}{:}~\PY{n}{TypeJG}~\PY{n}{Γ}~\PY{n}{A}\PY{o}{)}~\PY{o}{:}~\PY{n}{CtxJG}~\PY{n}{Γ}\\
    \mk{0}\PY{k}{with}~~\PY{n}{subs\PYZus{}wf}~\PY{n}{Δ}~\PY{n}{Γ}~\PY{n}{y}~\PY{o}{(}\PY{n}{H}~\PY{o}{:}~\PY{n}{SubsJG}~\PY{n}{Δ}~\PY{n}{Γ}~\PY{n}{y}\PY{o}{)}~\PY{o}{:}\\
    \mk{0}~~~~~~~~\PY{n}{CtxJG}~\PY{n}{Δ}~\PY{o}{/}\PY{o}{\PYZbs{}}~\PY{n}{CtxJG}~\PY{n}{Γ}\PY{o}{.}\\
    \mk{1}\PY{k+kn}{Proof}\PY{o}{.}\\
    \mk{2}\PY{o}{\PYZhy{}}~\PY{n+nb+bp}{exact}~\PY{o}{(}\PY{n}{type\PYZus{}wf}~\PY{n}{Γ}~\PY{n}{A}~\PY{n}{H}\PY{o}{)}\PY{o}{.}~~\PY{c}{(*}\PY{c}{~accepted~}\PY{c}{*)}\\
    \mk{3}\PY{o}{\PYZhy{}}~\PY{n+nb}{dependent~induction}~\PY{n}{H}\PY{o}{.}\\
    \mk{4}~~\PY{o}{+}~\PY{n+nb}{pose~proof}~\PY{o}{(}\PY{n}{type\PYZus{}wf}~\PY{n}{\PYZus{}}~\PY{n}{\PYZus{}}~\PY{n}{H}\PY{o}{)}~\PY{k}{as}~\PY{n}{HG}\PY{o}{.}\\
    \mk{5}~~~~\PY{n+nb}{split}\PY{o}{;}~\PY{o}{[}\PY{n+nb+bp}{exact}~\PY{o}{(}\PY{n}{ctx\PYZus{}ext}~\PY{n}{\PYZus{}}~\PY{n}{\PYZus{}}~\PY{n}{HG}~\PY{n}{H}\PY{o}{)}~\PY{o}{|}~\PY{n+nb+bp}{exact}~\PY{n}{HG}\PY{o}{]}\PY{o}{.}\\
    \mk{0}~~\PY{o}{+}~\PY{c}{(*}\PY{c}{~abridged~}\PY{c}{*)}\\
    \mk{6}\PY{k+kn}{Qed}\PY{o}{.}\\
    }
    \column{0.51\textwidth}
    \alectryonvscode\vspace*{-5pt}
    \begin{onlyenv}<1>
\begin{alectryon}{coq}{}
Proof. (* .unfold .no-in -.g#2.h{*} *)
\end{alectryon}
    \end{onlyenv}
    \begin{onlyenv}<2>
\begin{alectryon}{coq}{}
- (* .none *) exact (type_wf Γ A H).  (* accepted *) (* .none *)
\end{alectryon}
    \rocqmsg{This subproof is complete, but there are some unfocused goals.\\ Focus next goal with bullet -.}
    \end{onlyenv}
    \begin{onlyenv}<3>
\begin{alectryon}{coq}{}
- (* .none *) dependent induction H. (* .unfold .no-in -.g#2.h{*} *)
\end{alectryon}
    \end{onlyenv}
    \begin{onlyenv}<4>
\begin{alectryon}{coq}{}
+ (* .none *) pose proof (type_wf _ _ H) as HG. (* .unfold .no-in *)
\end{alectryon}
    \end{onlyenv}
    \begin{onlyenv}<5>
\begin{alectryon}{coq}{}
split; [exact (ctx_ext _ _ HG H) | exact HG]. (* .none *)
\end{alectryon}
    \rocqmsg{This subproof is complete, but there are some unfocused goals.\\ Focus next goal with bullet +.}
    \end{onlyenv}
\begin{alectryon}{coq}{none}
+ destruct IHSubsJG as [HD HG].
  split; [exact HD | exact (ctx_ext _ _ HG H0)].
\end{alectryon}
    \begin{onlyenv}<6>
\begin{alectryon}{coq}{}
Fail Qed. (* .none *)
\end{alectryon}
    \colorbox{red!8}{\parbox{.95\linewidth}{\ttfamily\scriptsize Error: Cannot guess decreasing argument of fix.}}
    \end{onlyenv}
  \end{columns}
\begin{alectryon}{coq}{none}
Abort.
End Bad.
\end{alectryon}
  \popup<7>{\begin{itemize}
      \item more freedom, less guidance: termination is checked only at \texttt{Qed}.
      \item the error says neither which lemma nor which line
    \end{itemize}}
\end{frame}

% --- [Guarded] -----------------------------------------------------
\begin{frame}[t]{Checking early: \texttt{Guarded.}}
  Same proof, with a checkpoint after the suspicious line.
\begin{alectryon}{coq}{none}
Module Grd.
Import Method2.
Set Printing Width 48.
Lemma type_wf Γ A (H : TypeJG Γ A) : CtxJG Γ
with  subs_wf Δ Γ y (H : SubsJG Δ Γ y) :
        CtxJG Δ /\ CtxJG Γ.
\end{alectryon}
  \vspace{1.5ex}
  \begin{columns}[t,totalwidth=1.06\textwidth]
    \column{0.55\textwidth}
    {\ttfamily\scriptsize\setlength{\parindent}{0pt}
    \mbox{}\\[-1pt]
    \mk{0}\PY{k+kn}{Lemma}~\PY{n+nf}{type\PYZus{}wf}~\PY{n+nv}{Γ}~\PY{n+nv}{A}~\PY{o}{(}\PY{n+nv}{H}~\PY{o}{:}~\PY{n}{TypeJG}~\PY{n}{Γ}~\PY{n}{A}\PY{o}{)}~\PY{o}{:}~\PY{n}{CtxJG}~\PY{n}{Γ}\\
    \mk{0}\PY{k}{with}~~\PY{n}{subs\PYZus{}wf}~\PY{n}{Δ}~\PY{n}{Γ}~\PY{n}{y}~\PY{o}{(}\PY{n}{H}~\PY{o}{:}~\PY{n}{SubsJG}~\PY{n}{Δ}~\PY{n}{Γ}~\PY{n}{y}\PY{o}{)}~\PY{o}{:}\\
    \mk{0}~~~~~~~~\PY{n}{CtxJG}~\PY{n}{Δ}~\PY{o}{/}\PY{o}{\PYZbs{}}~\PY{n}{CtxJG}~\PY{n}{Γ}\PY{o}{.}\\
    \mk{1}\PY{k+kn}{Proof}\PY{o}{.}\\
    \mk{2}\PY{o}{\PYZhy{}}~\PY{n+nb+bp}{exact}~\PY{o}{(}\PY{n}{type\PYZus{}wf}~\PY{n}{Γ}~\PY{n}{A}~\PY{n}{H}\PY{o}{)}\PY{o}{.}~~\PY{c}{(*}\PY{c}{~accepted~}\PY{c}{*)}\\
    \mk{3}~~\PY{k+kn}{Guarded}\PY{o}{.}\\
    \mk{0}\PY{o}{\PYZhy{}}~\PY{n+nb}{dependent~induction}~\PY{n}{H}\PY{o}{.}\\
    \mk{0}~~\PY{o}{+}~\PY{n+nb}{pose~proof}~\PY{o}{(}\PY{n}{type\PYZus{}wf}~\PY{n}{\PYZus{}}~\PY{n}{\PYZus{}}~\PY{n}{H}\PY{o}{)}~\PY{k}{as}~\PY{n}{HG}\PY{o}{.}\\
    \mk{0}~~~~\PY{n+nb}{split}\PY{o}{;}~\PY{o}{[}\PY{n+nb+bp}{exact}~\PY{o}{(}\PY{n}{ctx\PYZus{}ext}~\PY{n}{\PYZus{}}~\PY{n}{\PYZus{}}~\PY{n}{HG}~\PY{n}{H}\PY{o}{)}~\PY{o}{|}~\PY{n+nb+bp}{exact}~\PY{n}{HG}\PY{o}{]}\PY{o}{.}\\
    \mk{0}~~\PY{o}{+}~\PY{c}{(*}\PY{c}{~abridged~}\PY{c}{*)}\\
    \mk{0}\PY{k+kn}{Qed}\PY{o}{.}\\
    }
    \column{0.51\textwidth}
    \alectryonvscode\vspace*{-5pt}
    \begin{onlyenv}<1>
\begin{alectryon}{coq}{}
Proof. (* .unfold .no-in -.g#2.h{*} *)
\end{alectryon}
    \end{onlyenv}
    \begin{onlyenv}<2>
\begin{alectryon}{coq}{}
- (* .none *) exact (type_wf Γ A H).  (* accepted *) (* .none *)
\end{alectryon}
    \rocqmsg{This subproof is complete, but there are some unfocused goals.\\ Focus next goal with bullet -.}
    \end{onlyenv}
    \begin{onlyenv}<3>
\begin{alectryon}{coq}{}
Fail Guarded. (* .none *)
\end{alectryon}
    \colorbox{red!8}{\parbox{.95\linewidth}{\ttfamily\scriptsize Error: Recursive definition of type\_wf and subs\_wf is ill-formed.\\ As a mutual fixpoint decreasing on the 1st argument of type\_wf and 1st argument of subs\_wf: Recursive call to type\_wf has principal argument equal to "Γ" instead of a subterm of "Γ".\\ As a mutual fixpoint decreasing on the 1st argument of type\_wf and 2nd argument of subs\_wf: \dots\\ \dots}}\\{\tiny (error abridged: 12 combinations, 64 lines)}
    \end{onlyenv}
  \end{columns}
\begin{alectryon}{coq}{none}
Abort.
End Grd.
\end{alectryon}
  \popup<4>{\begin{itemize}
      \item \texttt{Guarded.}: check \emph{now}, as a checkpoint
      \item but it tries \textbf{every} combination of decreasing arguments:
            here $3 \times 4 = 12$ combinations, 64 lines of error
      \item exponential in the number of lemmas
    \end{itemize}}
\end{frame}

% --- [M1-whole] ----------------------------------------------------
\section{Looking back}
\begin{frame}[t]{Looking back: an induction principle, no termination check}
\begin{alectryon}{coq}{none}
Module Whole.
Import Method1.
Set Printing Width 48.
Definition wf_motive (f : Sort) : Prop :=
  match f with
  | TY Γ _   => J (CX Γ)
  | SB Δ Γ _ => J (CX Δ) /\ J (CX Γ)
  | TM Γ _ _ => J (CX Γ)
  | _        => True
  end.
Lemma wf_all f (H : J f) : wf_motive f.
\end{alectryon}
  \begin{columns}[t,totalwidth=1.06\textwidth]
    \column{0.55\textwidth}
    {\ttfamily\scriptsize\setlength{\parindent}{0pt}
    \mbox{}\\[-1pt]
    \mk{0}\PY{k+kn}{Definition}~\PY{n+nf}{wf\PYZus{}motive}~\PY{o}{(}\PY{n+nv}{f}~\PY{o}{:}~\PY{n}{Sort}\PY{o}{)}~\PY{o}{:}~\PY{k+kt}{Prop}~\PY{o}{:}\PY{o}{=}\\
    \mk{0}~~\PY{k}{match}~\PY{n}{f}~\PY{k}{with}\\
    \mk{0}~~\PY{o}{|}~\PY{n}{TY}~\PY{n}{Γ}~\PY{n}{\PYZus{}}~~~\PY{o}{=}\PY{o}{\PYZgt{}}~\PY{n}{J}~\PY{o}{(}\PY{n}{CX}~\PY{n}{Γ}\PY{o}{)}\\
    \mk{0}~~\PY{o}{|}~\PY{n}{SB}~\PY{n}{Δ}~\PY{n}{Γ}~\PY{n}{\PYZus{}}~\PY{o}{=}\PY{o}{\PYZgt{}}~\PY{n}{J}~\PY{o}{(}\PY{n}{CX}~\PY{n}{Δ}\PY{o}{)}~\PY{o}{/}\PY{o}{\PYZbs{}}~\PY{n}{J}~\PY{o}{(}\PY{n}{CX}~\PY{n}{Γ}\PY{o}{)}\\
    \mk{0}~~\PY{o}{|}~\PY{n}{TM}~\PY{n}{Γ}~\PY{n}{\PYZus{}}~\PY{n}{\PYZus{}}~\PY{o}{=}\PY{o}{\PYZgt{}}~\PY{n}{J}~\PY{o}{(}\PY{n}{CX}~\PY{n}{Γ}\PY{o}{)}\\
    \mk{0}~~\PY{o}{|}~\PY{n}{\PYZus{}}~~~~~~~~\PY{o}{=}\PY{o}{\PYZgt{}}~\PY{k+kt}{True}\\
    \mk{0}~~\PY{k}{end}\PY{o}{.}\\
    \mk{0}\PY{k+kn}{Lemma}~\PY{n+nf}{wf\PYZus{}all}~\PY{n+nv}{f}~\PY{o}{(}\PY{n+nv}{H}~\PY{o}{:}~\PY{n}{J}~\PY{n}{f}\PY{o}{)}~\PY{o}{:}~\PY{n}{wf\PYZus{}motive}~\PY{n}{f}\PY{o}{.}\\
    \mk{1}\PY{k+kn}{Proof}\PY{o}{.}\\
    \mk{2}~~\PY{n+nb}{induction}~\PY{n}{H}\PY{o}{;}~\PY{n+nb}{simpl}~\PY{k}{in}~\PY{o}{*}\PY{o}{.}\\
    \mk{3}~~\PY{o}{\PYZhy{}}~\PY{n+nb}{trivial}\PY{o}{.}\\
    \mk{4}~~\PY{o}{\PYZhy{}}\\
    \mk{5}~~~~\PY{n+nb+bp}{exact}~\PY{o}{(}\PY{n}{proj1}~\PY{n}{IHJ1}\PY{o}{)}\PY{o}{.}\\
    \mk{6}~~\PY{o}{\PYZhy{}}~\PY{n+nb+bp}{exact}~\PY{o}{(}\PY{n}{proj1}~\PY{n}{IHJ1}\PY{o}{)}\PY{o}{.}\\
    \mk{7}~~\PY{o}{\PYZhy{}}~\PY{n+nb}{split}\PY{o}{;}~\PY{o}{[}\PY{n+nb+bp}{exact}~\PY{o}{(}\PY{n}{ctx\PYZus{}ext}~\PY{n}{\PYZus{}}~\PY{n}{\PYZus{}}~\PY{n}{IHJ}~\PY{n}{H}\PY{o}{)}~\PY{o}{|}~\PY{n+nb+bp}{exact}~\PY{n}{IHJ}\PY{o}{]}\PY{o}{.}\\
    \mk{8}~~\PY{o}{\PYZhy{}}\\
    \mk{9}~~~~\PY{n+nb}{split}\PY{o}{;}~\PY{o}{[}\PY{n+nb+bp}{exact}~\PY{n}{IHJ3}~\PY{o}{|}~\PY{n+nb+bp}{exact}~\PY{o}{(}\PY{n}{ctx\PYZus{}ext}~\PY{n}{\PYZus{}}~\PY{n}{\PYZus{}}~\PY{n}{IHJ2}~\PY{n}{H0}\PY{o}{)}\PY{o}{]}\PY{o}{.}\\
    \mk{10}\PY{k+kn}{Qed}\PY{o}{.}\\
    }
    \column{0.51\textwidth}
    \alectryonvscode\vspace*{-5pt}
    \begin{onlyenv}<1>
\begin{alectryon}{coq}{}
Proof. (* .unfold .no-in *)
\end{alectryon}
    \end{onlyenv}
    \begin{onlyenv}<2>
\begin{alectryon}{coq}{}
induction H; simpl in *. (* .unfold .no-in -.g#2.h{*} -.g#3.h{*} -.g#4.h{*} -.g#5.h{*} *)
\end{alectryon}
    \end{onlyenv}
    \begin{onlyenv}<3>
\begin{alectryon}{coq}{}
- (* .none *) trivial. (* .none *)
\end{alectryon}
    \rocqmsg{This subproof is complete, but there are some unfocused goals.\\ Focus next goal with bullet -.}
    \end{onlyenv}
    \begin{onlyenv}<4>
\begin{alectryon}{coq}{}
- (* .unfold .no-in *)
\end{alectryon}
    \end{onlyenv}
    \begin{onlyenv}<5>
\begin{alectryon}{coq}{}
exact (proj1 IHJ1). (* .none *)
\end{alectryon}
    \rocqmsg{This subproof is complete, but there are some unfocused goals.\\ Focus next goal with bullet -.}
    \end{onlyenv}
    \begin{onlyenv}<6>
\begin{alectryon}{coq}{}
- (* .none *) exact (proj1 IHJ1). (* .none *)
\end{alectryon}
    \rocqmsg{This subproof is complete, but there are some unfocused goals.\\ Focus next goal with bullet -.}
    \end{onlyenv}
    \begin{onlyenv}<7>
\begin{alectryon}{coq}{}
- (* .none *) split; [exact (ctx_ext _ _ IHJ H) | exact IHJ]. (* .none *)
\end{alectryon}
    \rocqmsg{This subproof is complete, but there are some unfocused goals.\\ Focus next goal with bullet -.}
    \end{onlyenv}
    \begin{onlyenv}<8>
\begin{alectryon}{coq}{}
- (* .unfold .no-in *)
\end{alectryon}
    \end{onlyenv}
    \begin{onlyenv}<9>
\begin{alectryon}{coq}{}
split; [exact IHJ3 | exact (ctx_ext _ _ IHJ2 H0)]. (* .none *)
\end{alectryon}
    \rocqmsg{No more goals.}
    \end{onlyenv}
    \begin{onlyenv}<10>
\begin{alectryon}{coq}{}
Qed. (* .none *)
\end{alectryon}
    {\color{nagoyaGreen}\small$\checkmark$ accepted}
    \end{onlyenv}
  \end{columns}
\begin{alectryon}{coq}{none}
Lemma wf Γ A (H : J (TY Γ A)) : J (CX Γ).
Proof. exact (wf_all _ H). Qed.
End Whole.
\end{alectryon}
  \popup<11>{\begin{itemize}
      \item no \texttt{fix} in the proof: no termination check at \texttt{Qed}
      \item method 2 gets the same with \texttt{Combined Scheme}
      \item (c) done right: the induction comes \emph{before} the cases
    \end{itemize}}
\end{frame}

% --- [Takeaways] ---------------------------------------------------
\section{Conclusion}
\begin{frame}{Takeaways}
  \begin{itemize}
    \item A cycle of judgements: one inductive per judgement, tied with \texttt{with}
    \item A lemma on one judgement needs the others, in both methods
    \item Two ways to prove them together, available in both methods:
          \begin{itemize}
            \item (c) one induction principle on the \emph{whole} syntax, no termination check
            \item (d) \textbf{mutual lemmas} \texttt{Lemma \dots\ with \dots}: little known, mentioned once in the manual;
                  termination checked only at \texttt{Qed}, use \texttt{Guarded.} to check early;
                  \texttt{Fixpoint \dots\ \{struct H\}} seems to fix that (supplement)
          \end{itemize}
    \item Method 1 leans toward (c), method 2 toward (d)
    \item I recommend method 2: no sort, cycles are easy to add or remove,
          with \texttt{Scheme} both (c) and (d) are available, \dots
  \end{itemize}
  \vfill
  \centering
  {\Large Questions?}
\end{frame}

% =====================================================================
% Supplements (after the talk; not counted in the slide numbers)
% =====================================================================
\appendix
\section{Supplements}
% --- [Supp-Fix] ----------------------------------------------------
\begin{frame}[t]{Supplement: \texttt{Fixpoint \dots\ \{struct H\}} with no body}
  Also enters proof mode, but the decreasing argument is \textbf{given}.
\begin{alectryon}{coq}{none}
Module Fix.
Import Method2.
Set Printing Width 48.
Fixpoint type_wf Γ A (H : TypeJG Γ A)
        {struct H} : CtxJG Γ
with subs_wf Δ Γ y (H : SubsJG Δ Γ y)
        {struct H} :
        CtxJG Δ /\ CtxJG Γ.
\end{alectryon}
  \vspace{1.5ex}
  \begin{columns}[t,totalwidth=1.06\textwidth]
    \column{0.55\textwidth}
    {\ttfamily\scriptsize\setlength{\parindent}{0pt}
    \mbox{}\\[-1pt]
    \mk{0}\PY{k+kn}{Fixpoint}~\PY{n+nf}{type\PYZus{}wf}~\PY{n+nv}{Γ}~\PY{n+nv}{A}~\PY{o}{(}\PY{n+nv}{H}~\PY{o}{:}~\PY{n}{TypeJG}~\PY{n}{Γ}~\PY{n}{A}\PY{o}{)}\\
    \mk{0}~~~~~~~~\PY{o}{\PYZob{}}\PY{k}{struct}~\PY{n}{H}\PY{o}{\PYZcb{}}~\PY{o}{:}~\PY{n}{CtxJG}~\PY{n}{Γ}\\
    \mk{0}\PY{k}{with}~\PY{n}{subs\PYZus{}wf}~\PY{n}{Δ}~\PY{n}{Γ}~\PY{n}{y}~\PY{o}{(}\PY{n}{H}~\PY{o}{:}~\PY{n}{SubsJG}~\PY{n}{Δ}~\PY{n}{Γ}~\PY{n}{y}\PY{o}{)}\\
    \mk{0}~~~~~~~~\PY{o}{\PYZob{}}\PY{k}{struct}~\PY{n}{H}\PY{o}{\PYZcb{}}~\PY{o}{:}\\
    \mk{0}~~~~~~~~\PY{n}{CtxJG}~\PY{n}{Δ}~\PY{o}{/}\PY{o}{\PYZbs{}}~\PY{n}{CtxJG}~\PY{n}{Γ}\PY{o}{.}\\
    \mk{1}\PY{k+kn}{Proof}\PY{o}{.}\\
    \mk{2}\PY{o}{\PYZhy{}}~\PY{n+nb+bp}{exact}~\PY{o}{(}\PY{n}{type\PYZus{}wf}~\PY{n}{Γ}~\PY{n}{A}~\PY{n}{H}\PY{o}{)}\PY{o}{.}~~\PY{c}{(*}\PY{c}{~accepted~}\PY{c}{*)}\\
    \mk{3}~~\PY{k+kn}{Guarded}\PY{o}{.}\\
    \mk{4}\PY{o}{\PYZhy{}}~\PY{n+nb}{dependent~induction}~\PY{n}{H}\PY{o}{.}\\
    \mk{5}~~\PY{o}{+}~\PY{n+nb}{pose~proof}~\PY{o}{(}\PY{n}{type\PYZus{}wf}~\PY{n}{\PYZus{}}~\PY{n}{\PYZus{}}~\PY{n}{H}\PY{o}{)}~\PY{k}{as}~\PY{n}{HG}\PY{o}{.}\\
    \mk{6}~~~~\PY{n+nb}{split}\PY{o}{;}~\PY{o}{[}\PY{n+nb+bp}{exact}~\PY{o}{(}\PY{n}{ctx\PYZus{}ext}~\PY{n}{\PYZus{}}~\PY{n}{\PYZus{}}~\PY{n}{HG}~\PY{n}{H}\PY{o}{)}~\PY{o}{|}~\PY{n+nb+bp}{exact}~\PY{n}{HG}\PY{o}{]}\PY{o}{.}\\
    \mk{0}~~\PY{o}{+}~\PY{c}{(*}\PY{c}{~abridged~}\PY{c}{*)}\\
    \mk{7}\PY{k+kn}{Qed}\PY{o}{.}\\
    }
    \column{0.51\textwidth}
    \alectryonvscode\vspace*{-5pt}
    \begin{onlyenv}<1>
\begin{alectryon}{coq}{}
Proof. (* .unfold .no-in -.g#2.h{*} *)
\end{alectryon}
    \end{onlyenv}
    \begin{onlyenv}<2>
\begin{alectryon}{coq}{}
- (* .none *) exact (type_wf Γ A H).  (* accepted *) (* .none *)
\end{alectryon}
    \rocqmsg{This subproof is complete, but there are some unfocused goals.\\ Focus next goal with bullet -.}
    \end{onlyenv}
    \begin{onlyenv}<3>
\begin{alectryon}{coq}{}
Fail Guarded. (* .none *)
\end{alectryon}
    \colorbox{red!8}{\parbox{.95\linewidth}{\ttfamily\scriptsize Error: Recursive definition of type\_wf and subs\_wf is ill-formed.\\ As a mutual fixpoint decreasing on the 3rd argument of type\_wf and 4th argument of subs\_wf: Recursive call to type\_wf has principal argument equal to "H" instead of a subterm of "H".\\ \dots}}\\{\tiny (error abridged: 1 combination, 9 lines)}
    \end{onlyenv}
    \begin{onlyenv}<4>
\begin{alectryon}{coq}{}
- (* .none *) dependent induction H. (* .unfold .no-in -.g#2.h{*} *)
\end{alectryon}
    \end{onlyenv}
    \begin{onlyenv}<5>
\begin{alectryon}{coq}{}
+ (* .none *) pose proof (type_wf _ _ H) as HG. (* .unfold .no-in *)
\end{alectryon}
    \end{onlyenv}
    \begin{onlyenv}<6>
\begin{alectryon}{coq}{}
split; [exact (ctx_ext _ _ HG H) | exact HG]. (* .none *)
\end{alectryon}
    \rocqmsg{This subproof is complete, but there are some unfocused goals.\\ Focus next goal with bullet +.}
    \end{onlyenv}
\begin{alectryon}{coq}{none}
+ destruct IHSubsJG as [HD HG].
  split; [exact HD | exact (ctx_ext _ _ HG H0)].
\end{alectryon}
    \begin{onlyenv}<7>
\begin{alectryon}{coq}{}
Fail Qed. (* .none *)
\end{alectryon}
    \colorbox{red!8}{\parbox{.95\linewidth}{\ttfamily\scriptsize Error: Recursive definition of type\_wf is ill-formed.\\ \dots\\ Recursive call to type\_wf has principal argument equal to "H" instead of a subterm of "H".}}\\{\tiny (error abridged)}
    \end{onlyenv}
  \end{columns}
\begin{alectryon}{coq}{none}
Abort.
End Fix.
\end{alectryon}
  \popup<8>{\begin{itemize}
      \item one combination checked; the error names the lemma and the bad call
      \item \texttt{Guarded.} too: 1 combination, 9 lines (12 combinations, 64 lines without \texttt{struct})
      \item trade-off: you must know the decreasing argument \emph{before} the proof
    \end{itemize}}
\end{frame}

% --- [Supp-M1-ind] -------------------------------------------------
\begin{frame}{Supplement: method 1, induction on a judgement}
  A well-formed type lives in a well-formed context.
  \vspace{-1ex}
\begin{alectryon}{coq}{none}
Module Ind1.
Import Method1.
Set Printing Width 90.
\end{alectryon}
  \small
\begin{alectryon}{coq}{}
Lemma wf Γ A : J (TY Γ A) -> J (CX Γ). (* .in *)
Proof. (* .none *)
intros H; dependent induction H.
\end{alectryon}
\begin{alectryon}{coq}{none}
Abort.
End Ind1.
\end{alectryon}

  \vfill
  \begin{itemize}
    \item The IHs look like progress, but neither helps:\\
          \texttt{IHJ1} needs \texttt{SB \dots{} = TY \dots} (never true);
          \texttt{IHJ2} gives \texttt{J (CX Γ0)}, not \texttt{J (CX Γ)}
    \item What is missing: from \texttt{H : J (SB Γ Γ0 y)} to the goal \texttt{J (CX Γ)}
  \end{itemize}
\end{frame}

\end{document}
